% id: 69-13
% book: 베이직쎈 확률과 통계 · 04 조건부확률
% section: 유형 039 표를 이용한 조건부확률
% figure: none
% answer: 
% answer_source: 
% variant_level: 0 (원본 전사)
% 이 파일은 build.mjs 가 items.json 에서 만든다. 고칠 때는 items.json 을 고친다.
\smash{\raisebox{1.5mm}{\dmkichul{베이직쎈 확률과 통계 69쪽 13번}}}
아래는 어느 고등학교 1학년 학생 100명을 대상으로 등교할 때 자전거를 이용하는지 조사하여 나타낸 표이다. 조사 대상 학생 중 임의로 한 명을 택할 때, 다음을 구하시오.\par\noindent\hfill{\footnotesize (단위: 명)}\par\smallskip\dispeq{\begin{tabular}{|c|c|c|c|}\hline & 남학생 & 여학생 & 합계 \\ \hline 자전거 이용 & 24 & 16 & 40 \\ \hline 자전거 이용 안 함 & 30 & 30 & 60 \\ \hline 합계 & 54 & 46 & 100 \\ \hline\end{tabular}}\par\smallskip\noindent
\dmsubprob{1} 택한 학생이 자전거를 이용할 때, 그 학생이 남학생일 확률
\dmsubprob{2} 택한 학생이 남학생일 때, 그 학생이 자전거를 이용할 확률
